\subsection{Agentic Software Engineering and Spec-Driven AI-Native Lifecycles}
\label{subsec:rw1}

The integration of large language models (LLMs) into software construction has progressed rapidly from single-shot code completion to autonomous agents that plan, edit, execute, and repair across an entire repository. A first generation of techniques established that an LLM's output improves markedly when it is embedded in an iterative loop rather than emitted in one pass. \emph{Self-Refine} demonstrated that a single model can critique and rewrite its own output without additional training~\cite{madaan2023selfrefine}, while \emph{Reflexion} showed that verbalized feedback retained in an episodic memory buffer steers an agent toward better subsequent attempts~\cite{shinn2023reflexion}. For code specifically, \emph{Self-Debugging} taught models to inspect execution traces and explain their own faults~\cite{chen2024selfdebug}, and \emph{AlphaCodium} recast competitive-programming synthesis as a multi-stage ``flow'' in which the model reasons, generates, and then repairs candidate solutions against tests~\cite{ridnik2024alphacodium}. These approaches share a common architecture---generate, observe a signal, regenerate---that prefigures the synthesis loop this article formalizes; their limitation, however, is the nature of that signal. Each closes its loop on either model self-assessment or a finite, model-authored test suite, so the feedback is both circular and incomplete: the same model that wrote the code also adjudicates it, and passing the chosen examples says nothing about untested inputs.

A second strand scaled the loop to full software engineering tasks against real repositories. \emph{SWE-agent} showed that a carefully designed agent--computer interface lets an LLM navigate, edit, and test a codebase to resolve genuine GitHub issues~\cite{yang2024sweagent}, evaluated on the now-standard SWE-bench benchmark of real-world issues and their reference patches~\cite{jimenez2024swebench}. In parallel, multi-agent frameworks decompose development across specialized roles: \emph{MetaGPT} encodes standardized operating procedures across product-manager, architect, and engineer roles~\cite{hong2024metagpt}; \emph{ChatDev} organizes a virtual software company whose agents converse through a structured ``chat chain'' spanning design, coding, and testing~\cite{qian2024chatdev}; and \emph{AutoGen} provides a general substrate for orchestrating conversing agents and tool use~\cite{wu2024autogen}. Comprehensive surveys document the breadth of this turn from LLMs to LLM-based agents for software engineering~\cite{liu2024llmagentse}. Across these systems the dominant correctness oracle remains test execution---reference tests in the benchmark case, or model-generated tests otherwise---which inherits the classical incompleteness of testing: it samples the input space and demonstrates the presence of behaviour, never its absence over all inputs.

Most directly related is the emerging family of \emph{spec-driven} and \emph{AI-native} development lifecycles, which reorganise the process so that a specification, rather than code, becomes the primary human-authored artefact. GitHub's \emph{Spec Kit} formalizes a four-phase workflow---specify, plan, tasks, implement---anchored by a ``constitution'' of immutable project principles~\cite{githubspeckit}. Amazon's \emph{AI-Driven Development Lifecycle} (AI-DLC) restructures the lifecycle into inception, construction, and operations phases with built-in quality gates and human approvals, including a reverse-engineering step for brownfield codebases~\cite{awsaidlc}, and the \emph{Kiro} IDE places the requirements-to-design-to-tasks spec workflow at the centre of the editor~\cite{awskiro}. These efforts correctly identify that, when an agent can regenerate an implementation cheaply, the durable artefact is the specification, and broader position papers argue that LLMs will reshape the SDLC accordingly~\cite{ozkaya2023llmse}. Their decisive shortcoming, however, is that the specification is expressed in \emph{natural language}. Consequently the specification-to-code step is a \emph{compilation without a checker}: there is no mechanical relation guaranteeing that the generated implementation satisfies the stated intent, the loop remains open (generate, then review by eye), and where verification appears at all it is the circular, self-grading test execution noted above. The natural-language artefact is also too imprecise to support automated reasoning over change: these lifecycles offer no way to compute, without re-running anything, which dependents a modification might break.

This article addresses precisely that missing verification spine. In contrast to the surveyed work, it proposes a \emph{machine-checkable formal-contract layer} as the pivot between ambiguous human intent and precise machine behaviour, in which the contract serves the dual role of input to synthesis and oracle for verification. Verification is performed by OpenJML Extended Static Checking over \emph{all} inputs rather than a sampled suite, closing the open loop with sound counterexamples instead of failing test indices, and an \emph{independent-authority constraint} forbids the agent that writes the code from authoring the contract it is judged against, eliminating the self-grading circularity endemic to the prior loops. Because formal contracts compose, the approach further treats a callee-specification change as a recomputation of caller proof obligations---turning ``which call sites does this upgrade break?'' into a first-class, run-free analysis that the natural-language lifecycles cannot express. Finally, recognising that most software is brownfield, the article contributes a multi-source characterisation phase that triangulates code, issue tickets, and documentation to recover a provenance-tagged candidate contract and to surface code-versus-intent disagreements as candidate bugs, going beyond the simple reverse-engineering step offered by AI-DLC.
